\documentclass[11pt,a4paper]{article}
\usepackage[T1]{fontenc}
\usepackage[utf8]{inputenc}
\usepackage{lmodern}
\usepackage{amsmath,amssymb,amsthm,mathtools}
\usepackage{esint,graphicx}
\usepackage[a4paper,margin=22mm,headheight=14pt]{geometry}
\usepackage{microtype}
\usepackage{booktabs,longtable,array,enumitem}
\newcolumntype{L}[1]{>{\raggedright\arraybackslash}p{#1}}
\usepackage[dvipsnames]{xcolor}
\usepackage{fancyhdr}
\usepackage{hyperref}

\newtheorem{theorem}{Theorem}[section]
\newtheorem{lemma}[theorem]{Lemma}
\newtheorem{proposition}[theorem]{Proposition}
\newtheorem{corollary}[theorem]{Corollary}
\theoremstyle{remark}
\newtheorem{remark}[theorem]{Remark}
\newtheorem{definition}[theorem]{Definition}
\newtheorem{claim}[theorem]{Claim}
\setlist{nosep,leftmargin=1.5em}
\setlength{\emergencystretch}{2em}
\pagestyle{fancy}
\fancyhf{}
\fancyhead[L]{\small AIM 690}
\fancyhead[R]{\small Revised 7 November 2026}
\fancyfoot[C]{\thepage}
\renewcommand{\headrulewidth}{0.3pt}
\hypersetup{colorlinks=true,linkcolor=MidnightBlue,urlcolor=MidnightBlue,citecolor=MidnightBlue,pdftitle={AIM 690},pdfauthor={Matthew J. Colbrook; results found by AI}}
\begin{document}
\title{Every logarithmic Smirnov function extends to the Hilbert polydisc}
\author{Matthew J. Colbrook\thanks{The solutions and partial results in this document were found by AI. The supplied research and the separate repository review are AI work; human peer review and formal verification are not claimed.}\\[4pt]
\small Department of Applied Mathematics and Theoretical Physics\\
\small University of Cambridge, Cambridge CB3 0WA, United Kingdom\\
\small\href{mailto:m.colbrook@damtp.cam.ac.uk}{m.colbrook@damtp.cam.ac.uk}}
\date{7 November 2026}
\maketitle
\begin{abstract}
Every radial Nevanlinna function on the $\ell^1$ polydisc extends uniquely and Fr\'echet holomorphically to the Hilbert polydisc. The latter is exactly the domain of continuous polynomial evaluation for the logarithmic metric. The conclusion concerns exactly the catalogue target; no historical priority claim is made.
\end{abstract}
\noindent\textbf{Catalogue scope.} AIM 690; complete resolution. \href{https://github.com/MColbrook/AIM/blob/1e564633494c109edfd3ec46da1366adcff925a5/problems/690-a-smirnov-function-without-a-hilbert-domain-holomorphic-extension.md}{Canonical target at revision \texttt{1e56463}}.


\section{Statement and argument}
\label{sec:690}

For $q=1,2$, write
$D_q^\infty=\ell^q(\mathbb C)\cap\mathbb D^\infty$, with its
$\ell^q$ topology.  The space $N_*(\mathbb T^\infty)$ is the closure
of analytic polynomials depending on finitely many coordinates for
the metric
\[
 d_0(f,g)=\int_{\mathbb T^\infty}\log(1+|f-g|)\,dm.
\]
The answer to the existence question is negative: every element of
$N_*(\mathbb T^\infty)$ has a Fr\'echet-holomorphic extension to
$D_2^\infty$, agreeing with its canonical representative on
$D_1^\infty$.

The key is a dimension-independent estimate for the logarithmic metric.
We construct a bounded averaging density from analytic discs; the
ordinary product Poisson density itself need not be bounded uniformly
in the dimension.  No infinite-dimensional change of variables is
needed.

Write $m_N$ for normalized Haar measure on $\mathbb T^N$.
For $0<r<1$ and $M\ge0$, set
\begin{equation}
 C(r,M)=\exp\left(
 2M+\left[1+\frac1{1-r^2}\right]M^2
 \right).
 \label{eq:690-1}
\end{equation}

\begin{lemma}[Dimension-independent logarithmic evaluation]
For every $N\ge1$, every $z\in\mathbb D^N$ with
$\max_j|z_j|\le r$ and $\sum_j|z_j|^2\le M^2$, and every analytic
polynomial $P$ in $N$ variables,
\begin{equation}
 \log(1+|P(z)|)
 \le C(r,M)\int_{\mathbb T^N}\log(1+|P(w)|)\,dm_N(w).
 \label{eq:690-2}
\end{equation}
\end{lemma}

\begin{proof}
Define the coordinatewise disc automorphism
\[
 T_z(v)_j=\frac{z_j+v_j}{1+\overline{z_j}v_j}.
\]
The pushforward $(T_z)_*m_N$ has the product Poisson density
\[
 P_z(w)=\prod_{j=1}^N
 \frac{1-|z_j|^2}{|1-\overline{z_j}w_j|^2}.
\]
In particular,
\begin{equation}
 P_z(T_z(v))
 =\prod_{j=1}^N
 \frac{|1+\overline{z_j}v_j|^2}{1-|z_j|^2}.
 \label{eq:690-3}
\end{equation}
For $v\in\mathbb T^N$, put
\begin{equation}
 S_z(v)=\sum_{j=1}^N\overline{z_j}v_j,\qquad
 Z_z=\int_{\mathbb T^N}e^{-2|S_z(v)|}\,dm_N(v),\qquad
 h_z(v)=\frac{e^{-2|S_z(v)|}}{Z_z}.
 \label{eq:690-4}
\end{equation}
The density $h_z$ is continuous and positive, and $h_zm_N$ is a
probability measure.  Crucially,
\begin{equation}
 h_z(\lambda v)=h_z(v)\qquad(|\lambda|=1).
 \label{eq:690-5}
\end{equation}

For each fixed $v\in\mathbb T^N$, the holomorphic disc
$\lambda\mapsto T_z(\lambda v)$ has center $z$ and extends
holomorphically past the closed unit disc.  The function
$\log(1+|P(T_z(\lambda v))|)$ is subharmonic, so
\[
 \log(1+|P(z)|)
 \le\int_{\mathbb T}
 \log(1+|P(T_z(\lambda v))|)\,dm_1(\lambda).
\]
Integrating against $h_z(v)dm_N(v)$ and using \eqref{eq:690-5} gives
\begin{equation}
 \log(1+|P(z)|)
 \le\int_{\mathbb T^N}
 \log(1+|P(T_z(v))|)h_z(v)\,dm_N(v).
 \label{eq:690-6}
\end{equation}
After $w=T_z(v)$, the probability measure in this integral has density
\[
 g_z(w)=h_z(T_z^{-1}(w))P_z(w).
\]
We now bound this density independently of $N$.
For $|u|<1$, the elementary inequality $\log(1+t)\le t$ gives
\[
 \log|1+u|^2
 =\log(1+2\operatorname{Re}u+|u|^2)
 \le2\operatorname{Re}u+|u|^2.
\]
Set
\[
 A_z=\|z\|_2^2-\sum_j\log(1-|z_j|^2).
\]
Equation \eqref{eq:690-3} then implies
\[
 P_z(T_z(v))\le\exp(2|S_z(v)|+A_z).
\]
Consequently
\begin{equation}
 g_z(T_z(v))
 =h_z(v)P_z(T_z(v))\le\frac{e^{A_z}}{Z_z}.
 \label{eq:690-7}
\end{equation}
Orthogonality of the coordinate functions and Cauchy--Schwarz give
\[
 \int_{\mathbb T^N}|S_z(v)|\,dm_N(v)
 \le\left(\sum_j|z_j|^2\right)^{1/2}.
\]
Jensen's inequality therefore yields
\begin{equation}
 Z_z\ge\exp\left(-2\int|S_z|\,dm_N\right)
 \ge e^{-2\|z\|_2}.
 \label{eq:690-8}
\end{equation}
Finally $-\log(1-t)\le t/(1-r^2)$ for $0\le t\le r^2$, so
\begin{equation}
 \frac{e^{A_z}}{Z_z}
 \le\exp\left(2\|z\|_2+
 \left[1+\frac1{1-r^2}\right]\|z\|_2^2\right)
 \le C(r,M).
 \label{eq:690-9}
\end{equation}
Combining \eqref{eq:690-6}--\eqref{eq:690-9} proves \eqref{eq:690-2}.
\end{proof}

\begin{theorem}
Every $f\in N_*(\mathbb T^\infty)$ admits a unique
Fr\'echet-holomorphic function $G:D_2^\infty\to\mathbb C$ whose
restriction to $D_1^\infty$ is the canonical representative of $f$.
\end{theorem}

\begin{proof}
By the definition of $N_*$, there are analytic polynomials $P_k$, each
depending on finitely many coordinates, such that $d_0(P_k,f)\to0$.
In particular $(P_k)$ is $d_0$-Cauchy.
For $0<r<1$ and $M<\infty$, write
\[
 V_{r,M}=\{\zeta\in\ell^2:
 \|\zeta\|_\infty\le r,\ \|\zeta\|_2\le M\}.
\]
Apply the lemma to $P_k-P_l$ using the finite set of its coordinates.
Projection onto this set does not increase either norm of $\zeta$,
and the finite-dimensional Haar integral is exactly the
infinite-dimensional integral for this cylindrical polynomial.  Thus
\begin{equation}
 \sup_{\zeta\in V_{r,M}}|P_k(\zeta)-P_l(\zeta)|
 \le\exp\bigl(C(r,M)d_0(P_k,P_l)\bigr)-1
 \longrightarrow0.
 \label{eq:690-10}
\end{equation}
There is therefore a well-defined limit $G$ on $D_2^\infty$, with
uniform convergence on every $V_{r,M}$.

The set $D_2^\infty$ is open in $\ell^2$.  Indeed, every
$\zeta\in D_2^\infty$ has $\|\zeta\|_\infty<1$: its coordinates
tend to zero, while each of the finitely many remaining coordinates
has modulus less than one.  Every such point has an open Hilbert-space
neighborhood contained in some $V_{r,M}$.  The $P_k$ are continuous
Fr\'echet-holomorphic polynomials on $\ell^2$; their locally uniform
limit is Fr\'echet holomorphic on $D_2^\infty$.

On $D_1^\infty$, the canonical representative is the pointwise limit
of the same $d_0$-approximating polynomials, by the construction in
\cite[Lemma 2.24 and Theorem 2.27]{GuoNiZhou2025}.
Thus $G$ has the required restriction.  Estimate \eqref{eq:690-10} also proves
independence of the polynomial approximation.  Uniqueness follows
from continuity and the density of $D_1^\infty$ in $D_2^\infty$,
witnessed by finite-support truncations.
\end{proof}

\begin{corollary}[Extension for the full radial Nevanlinna class]
Let $F$ be holomorphic on $D_1^\infty$ and suppose
\begin{equation}
 K:=\sup_{0<s<1}\int_{\mathbb T^\infty}
 \log^+\left|F(sw_1,s^2w_2,s^3w_3,\ldots)\right|\,dm(w)<\infty.
 \label{eq:690-11}
\end{equation}
Then $F$ extends uniquely to a Fr\'echet-holomorphic function on
$D_2^\infty$.  Thus the conclusion holds for the full class
$N(D_1^\infty)$ defined by \cite[Section 2.1, equation (2.3)]{GuoNiZhou2025},
as well as for its Smirnov subclass.
\end{corollary}

\begin{proof}
For $0<s<1$, let $D_s z=(s^jz_j)_{j\ge1}$ and set
$F_s=F\circ D_s$.  The map $D_s:\ell^2\to\ell^1$ is bounded,
so $F_s$ is holomorphic on $D_2^\infty$.
Moreover, the weighted closed polydisc
\[
 K_s=\prod_{j\ge1}s^j\overline{\mathbb D}
\]
is compact in $\ell^1$ and contained in $D_1^\infty$.
Uniform continuity of $F$ on $K_s$ shows that its finite-coordinate
truncations approximate $F_s$ uniformly on the closed coordinate
polydisc.  Each truncation is holomorphic on a neighborhood of its
finite-dimensional closed polydisc, and hence is uniformly
approximable there by analytic polynomials.  Thus $F_s$ belongs to the
infinite polydisc algebra, as also recorded immediately before the
proof of \cite[Proposition 2.3]{GuoNiZhou2025}.

Passing to uniform polynomial limits in \eqref{eq:690-2}, and using
$\log(1+x)\le\log2+\log^+x$, gives
\begin{equation}
 \log(1+|F_s(\zeta)|)
 \le C(r,M)(K+\log2)
 \qquad(\zeta\in V_{r,M},\ 0<s<1).
 \label{eq:690-12}
\end{equation}
The family $(F_s)_{0<s<1}$ is therefore locally uniformly bounded in
the Hilbert-space topology on $D_2^\infty$.

Fix a point $\zeta\in D_2^\infty$ and a Hilbert ball about it whose
doubled ball lies in some $V_{r,M}$.  Cauchy's estimate on complex
lines gives a uniform bound for the operator norms of $DF_s$ on the
smaller ball.  Consequently the $F_s$ are equi-Lipschitz there.
At every finite-support point $a$, $D_sa\to a$ in $\ell^1$ as
$s\uparrow1$, so $F_s(a)\to F(a)$.
Finite-support points are dense in $\ell^2$; comparing two functions
$F_s,F_t$ at a nearby such point and using the common Lipschitz bound
proves that $F_s(\zeta)$ is Cauchy as $s\uparrow1$.
Denote its limit by $G(\zeta)$.

The limit $G$ is locally bounded by \eqref{eq:690-12}.  On each compact disc
in any complex line in $D_2^\infty$, the common local Lipschitz bounds
and pointwise convergence imply uniform convergence, by a finite
net argument.  Hence the restriction of $G$ to every such line is
holomorphic.  Local boundedness and this G\^ateaux holomorphy imply
Fr\'echet holomorphy.  This argument does not require a
finite-dimensional Montel theorem on a noncompact Hilbert ball.

Finally, for every $\zeta\in D_1^\infty$, dominated convergence gives
$\|D_s\zeta-\zeta\|_1\to0$.  Therefore
$F_s(\zeta)\to F(\zeta)$, proving that $G$ extends $F$.
Density again proves uniqueness.
\end{proof}

The uniform bound \eqref{eq:690-10} also excludes polynomials with arbitrarily
small logarithmic metric but large evaluation at a fixed Hilbert-polydisc
point.  The averaging measure is correlated across coordinates through
their common rotation; this is why unboundedness of the usual product
Poisson density does not provide such a construction.

\begin{corollary}[Sharpness of the Hilbert threshold]
Among coordinate points $z\in\mathbb D^\infty$, polynomial evaluation
$P\mapsto P(z)$ is continuous for the logarithmic metric if and only
if $z\in\ell^2$.
\end{corollary}

\begin{proof}
Sufficiency follows from \eqref{eq:690-2}.  Conversely, let
$z\in\mathbb D^\infty\setminus\ell^2$ and put
$s_N^2=\sum_{j=1}^N|z_j|^2$.  For all sufficiently large $N$, define
\[
 Q_N(w)=\frac{\sum_{j=1}^N\overline{z_j}w_j}{s_N^2}.
\]
Then $Q_N(z)=1$, whereas orthogonality and Cauchy--Schwarz give
\[
 d_0(Q_N,0)\le\int_{\mathbb T^N}|Q_N|\,dm_N
 \le\left(\int_{\mathbb T^N}|Q_N|^2\,dm_N\right)^{1/2}
 =s_N^{-1}\longrightarrow0.
\]
This contradicts continuity of evaluation at such a point.
\end{proof}

\begingroup\small
\begin{thebibliography}{99}
\bibitem{GuoNiZhou2025} K. Guo, J. Ni and Q. Zhou, \emph{Nevanlinna class, Dirichlet series and Szeg\H{o}'s problem}, Ann. Inst. Fourier \textbf{75} (2025), 1603--1647. \url{https://doi.org/10.5802/aif.3696}. Numbered passages refer to \url{https://arxiv.org/abs/2201.01993v2}.
\end{thebibliography}
\endgroup
\end{document}
